% ============================================================
% resume.tex — 中文简历 LaTeX 模板（xelatex 编译）
% 编译命令: xelatex resume.tex
%
% 所有可调参数集中在"参数区"，修改数值即可。
% 内容区在文件后半部分，填入你自己的信息。
% ============================================================

% ╔══════════════════════════════════════════════════════════╗
% ║                    文档类与包                            ║
% ╚══════════════════════════════════════════════════════════╝

\documentclass[11pt,a4paper]{ctexart}

% --- 页面布局 ---
\usepackage[top=1.5cm, bottom=1.2cm, left=1.5cm, right=1.5cm]{geometry}

% --- 字体 ---
\usepackage{fontspec}
% 首选 Microsoft YaHei，若缺失则回退到 FandolSong
\IfFontExistsTF{Microsoft YaHei}{%
  \setmainfont{Microsoft YaHei}
  \setsansfont{Microsoft YaHei}
}{%
  \setmainfont{FandolSong}[BoldFont=FandolHei]
  \setsansfont{FandolHei}
}

% --- 颜色 ---
\usepackage[svgnames]{xcolor}
\definecolor{accent}{HTML}{1a5276}        % 深蓝色，主色调
\definecolor{muted}{HTML}{666666}         % 灰色，次要文字
\definecolor{lightgray}{HTML}{eeeeee}     % 浅灰，分割线等

% --- 表格 ---
\usepackage{tabularx}
\usepackage{array}

% --- 列表（紧凑排版） ---
\usepackage{enumitem}
\setlist{noitemsep, topsep=0pt, partopsep=0pt, parsep=0pt, itemsep=0pt}

% --- 段落间距 ---
\usepackage{parskip}
\setlength{\parskip}{0pt}
\setlength{\parindent}{0pt}

% --- 彩色盒子（板块标题色条） ---
\usepackage{tcolorbox}
\tcbuselibrary{skins, breakable}

% --- 图片（头像占位） ---
\usepackage{graphicx}

% --- 绘图 ---
\usepackage{tikz}

% --- 超链接 ---
\usepackage[hidelinks]{hyperref}

% --- 行距 ---
\linespread{1.05}
\setlength{\intextsep}{0pt}
\setlength{\textfloatsep}{0pt}

% ╔══════════════════════════════════════════════════════════╗
% ║                    参数区（改这里）                      ║
% ╚══════════════════════════════════════════════════════════╝

% 字号
\newcommand{\sizeName}{\fontsize{22pt}{26pt}\selectfont}          % 姓名
\newcommand{\sizeRole}{\fontsize{12pt}{16pt}\selectfont}           % 职位
\newcommand{\sizeSection}{\fontsize{12pt}{14pt}\selectfont}        % 板块标题
\newcommand{\sizeBody}{\fontsize{9.5pt}{13pt}\selectfont}          % 正文
\newcommand{\sizeSmall}{\fontsize{8.5pt}{11pt}\selectfont}         % 小字（日期等）

% 间距
\newcommand{\gapBeforeSection}{10pt}   % 板块前间距
\newcommand{\gapAfterSection}{3pt}     % 板块后间距
\newcommand{\gapBeforeEntry}{6pt}      % 条目前间距
\newcommand{\gapAfterEntry}{2pt}       % 条目后间距

% ╔══════════════════════════════════════════════════════════╗
% ║                 自定义组件（不用动）                     ║
% ╚══════════════════════════════════════════════════════════╝

% ---------- 板块标题：左侧色条 ----------
\newcommand{\sectionTitle}[1]{%
  \vspace{\gapBeforeSection}
  \noindent
  \begin{tcolorbox}[
    enhanced,
    boxrule=0pt,
    leftrule=3pt,
    rightrule=0pt,
    toprule=0pt,
    bottomrule=0pt,
    colback=white,
    colframe=accent,
    arc=0pt,
    outer arc=0pt,
    boxsep=0pt,
    left=6pt,
    right=0pt,
    top=0pt,
    bottom=1pt,
    before skip=0pt,
    after skip=0pt,
    width=\textwidth,
    nobeforeafter,
  ]
  {\sizeSection\bfseries\color{accent}#1}
  \end{tcolorbox}
  \vspace{\gapAfterSection}
}

% ---------- 条目：标题行（机构 · 职位  |  日期） ----------
\newcommand{\entryHeader}[3]{%
  \vspace{\gapBeforeEntry}
  \noindent
  \begin{tabularx}{\textwidth}{@{}>{\bfseries\sizeBody}X@{\hspace{6pt}}>{\sizeSmall}r@{}}
    #1 \hspace{2pt}\textperiodcentered\hspace{2pt} #2 & #3
  \end{tabularx}
}

% ---------- 条目：描述要点 ----------
\newcommand{\entryBody}[1]{%
  \vspace{1pt}
  \begin{itemize}[leftmargin=12pt, label=--, itemsep=0pt, topsep=0pt]
    #1
  \end{itemize}
  \vspace{\gapAfterEntry}
}

% ---------- 技能组 ----------
\newcommand{\skillGroup}[2]{%
  \noindent
  {\bfseries\color{accent}\sizeBody#1：}\hspace{2pt}
  {\sizeBody#2}
  \par\vspace{1pt}
}

% ---------- 项目 ----------
\newcommand{\projectEntry}[3]{%
  \vspace{\gapBeforeEntry}
  \noindent
  {\bfseries\sizeBody#1}\hspace{4pt}{\sizeSmall\color{muted}#2}
  \vspace{1pt}
  \begin{itemize}[leftmargin=12pt, label=--, itemsep=0pt, topsep=0pt]
    #3
  \end{itemize}
  \vspace{\gapAfterEntry}
}

% ---------- 获奖条目 ----------
\newcommand{\awardEntry}[3]{%
  \noindent
  \begin{tabularx}{\textwidth}{@{}>{\sizeSmall\color{muted}}p{4.2cm}@{\hspace{4pt}}>{\bfseries\sizeBody}X@{\hspace{6pt}}>{\sizeSmall}r@{}}
    #1 & #2 & #3
  \end{tabularx}
  \par\vspace{1pt}
}

% ╔══════════════════════════════════════════════════════════╗
% ║                    正文开始                              ║
% ╚══════════════════════════════════════════════════════════╝

\begin{document}
\color{black!85}  % 正文用 85% 黑，视觉更柔和

\sizeBody

% ════════════════════════════════════════════════════════════
% 头部：头像 + 个人信息
% ════════════════════════════════════════════════════════════
\noindent
\begin{minipage}[t]{2.8cm}
  \vspace{0pt}
  % --- 头像占位框（替换 \includegraphics 为你的照片路径）---
  \begin{tikzpicture}
    \fill[accent!10] (0,0) rectangle (2.8,3.6);
    \node[text width=2.8cm, align=center, text=accent!50] at (1.4,1.8) {\small 照片\\\small 2.8×3.6cm};
  \end{tikzpicture}
\end{minipage}
\hfill
\begin{minipage}[t]{\dimexpr\textwidth-3.1cm\relax}
  \vspace{0pt}
  % 姓名
  \noindent{\sizeName\bfseries\color{accent}周康禧}

  \vspace{2pt}
  % 职位
  \noindent{\sizeRole\color{muted}C++量化开发实习生}

  \vspace{4pt}
  % 联系信息表格
  \noindent\sizeSmall
  \begin{tabularx}{\linewidth}{@{}>{\color{muted}}lX@{}}
    电话   & 15071352747 \\
    邮箱   & 2067545015@qq.com \\
    GitHub & github.com/yjwyyds1 \\
    地址   & 武汉市 \\
  \end{tabularx}
\end{minipage}

% ════════════════════════════════════════════════════════════
% 教育经历
% ════════════════════════════════════════════════════════════
\sectionTitle{教育经历}

\entryHeader{黄冈师范学院}{大数据科学与大数据技术 · 本科}{2023 -- 2027}
\entryBody{
  \item 主修课程：数据结构与算法、数据库技术、Linux操作系统基础、统计学等。
}

% ════════════════════════════════════════════════════════════
% 获奖经历
% ════════════════════════════════════════════════════════════
\sectionTitle{获奖经历}

\awardEntry{2025/11}{「华为杯」第50届ACM-ICPC国际大学生程序设计竞赛区域赛武汉站}{区域赛铜奖}
\awardEntry{2025/06}{CCPC中国大学生程序设计竞赛-全国邀请赛（郑州）}{邀请赛银奖}
\awardEntry{2026/05}{2026年ACM-ICPC国际大学生程序设计竞赛全国邀请赛（武汉）暨湖北省赛}{邀请赛银奖/省赛银奖}
\awardEntry{2026/06}{第十七届蓝桥杯全国软件和信息技术专业人才大赛 C/C++程序设计}{国赛二等奖}
\awardEntry{2026/05}{第十七届蓝桥杯全国软件和信息技术专业人才大赛 Python程序设计}{省赛一等奖}

% ════════════════════════════════════════════════════════════
% 专业技能
% ════════════════════════════════════════════════════════════
\sectionTitle{专业技能}

\skillGroup{语言}{C/C++ \quad Python \quad SQL}
\skillGroup{框架}{Spring Boot \quad Kafka \quad Redis \quad Flink \quad React \quad Node.js}
\skillGroup{基础设施}{Docker \quad Kubernetes \quad AWS \quad Terraform \quad MySQL \quad TiDB}
\skillGroup{工具}{Git \quad CI/CD \quad 微服务治理 \quad 性能调优 \quad 系统设计}

% ════════════════════════════════════════════════════════════
% 实习经历
% ════════════════════════════════════════════════════════════
\sectionTitle{实习经历}

\entryHeader{南京量化 · 数据分析与算法研发部}{C++开发实习生}{2026.04 -- 至今}
\entryBody{
  \item 负责抖音直播营收系统的微服务架构设计与性能优化，支撑日均 500 万+ 付费用户的高并发交易场景。
  \item 主导数据库读写分离与缓存层重构，将核心交易链路 P99 延迟从 380ms 降低至 85ms，可用性提升至 99.99\%。
  \item 带领 6 人后端团队，建立 Code Review 规范与线上事故分级响应机制，团队人效提升 40\%。
}

% ════════════════════════════════════════════════════════════
% 项目经历
% ════════════════════════════════════════════════════════════
\sectionTitle{项目经历}

\projectEntry{XAnalyzer --- 开源 Java 性能分析工具}{github.com/zhangmingyuan/xanalyzer}{
  \item 基于字节码插桩的 Java 应用性能分析 CLI 工具，支持热点方法检测、内存分配追踪与线程死锁诊断。GitHub 3.2k Stars。
}

\projectEntry{QuickDeploy --- 轻量级 CI/CD 平台}{quickdeploy.dev}{
  \item 基于 Go + Vue.js 构建的轻量级持续部署平台，支持 Docker 单机与 Kubernetes 集群两种模式。SaaS 月活团队 300+。
}

\end{document}
